\section{从 ODE 到 PDE：未知量从一条曲线变成一场}
\label{sec:15-Differential-Equations-and-Dynamical-Systems/03-PDE-Introduction/01}

一根铁杆，左端压在冰上，右端伸进火里。问它十秒钟后杆上的温度是多少——这个问题问不出一个数来。杆上每一点都有自己的温度，每一点的温度又都在变。ODE 求解器习惯交出的东西是一条曲线 \(x(t)\)，每个时刻对应一个数；这里要交出的是一张曲面 \(u(t,x)\)，每个时刻对应一整条剖面。维数从一根轴涨到两根轴，随之而来的麻烦不只是画图变难。

\subsection{把杆切成 N 段，得到的是 N 个绞在一起的 ODE}

先试着用手里现成的工具硬撑。在杆上取 \(N\) 个等距点 \(x_1,\dots,x_N\)，记 \(u_i(t)=u(t,x_i)\)，于是有了 \(N\) 个关于时间的未知函数。每个 \(u_i\) 的变化率由谁决定？由流进流出这一小段的热量决定，而热流正比于温度的斜率，斜率又要用邻居的值去算：\(\partial_x u\approx(u_{i+1}-u_i)/\Delta x\)。第 \(i\) 个方程的右端写不完，除非把 \(u_{i-1}\) 和 \(u_{i+1}\) 也写进去。

这 \(N\) 个 ODE 因此是绞在一起的，谁也不能单独求。加密网格能让逼近更准，代价是方程规模同步膨胀。而你要描述的物理并不在某个具体的 \(N\) 上，它在 \(N\to\infty\) 的极限里：差商收敛成导数，方程组收敛成一个方程，未知量收敛成一个定义在整段杆上的函数。到了那一步，语言必须换。

\begin{center}
\begin{tikzpicture}[>=stealth]
  \begin{scope}
    \draw[->] (0,0) -- (3.5,0) node[right] {\small \(t\)};
    \draw[->] (0,0) -- (0,2.5) node[above] {\small \(x\)};
    \draw[thick] plot[domain=0:3.2,samples=60,variable=\t]
      ({\t},{2.0-1.5*exp(-0.9*\t)});
    \draw[dashed] (1.6,0) -- (1.6,1.645);
    \fill (1.6,1.645) circle (1.8pt);
    \node at (1.75,2.35) {\small 一个时刻交出一个数};
  \end{scope}

  \begin{scope}[shift={(5.4,0)}]
    \draw[->] (0,0) -- (3.5,0) node[right] {\small \(x\)};
    \draw[->] (0,0) -- (0,2.5) node[above] {\small \(u\)};
    \draw[thick] plot[smooth] coordinates
      {(0,0.05) (0.5,0.12) (1.0,0.60) (1.3,1.55) (1.5,1.95)
       (1.7,1.55) (2.0,0.60) (2.5,0.12) (3.0,0.05)};
    \draw[dashed] plot[smooth] coordinates
      {(0,0.10) (0.5,0.32) (1.0,0.88) (1.5,1.25) (2.0,0.88)
       (2.5,0.32) (3.0,0.10)};
    \draw[dotted,thick] plot[smooth] coordinates
      {(0,0.16) (0.75,0.46) (1.5,0.74) (2.25,0.46) (3.0,0.16)};
    \foreach \x/\y in {0.75/0.28, 1.5/1.95, 2.25/0.28}{
      \draw[gray] (\x,0) -- (\x,\y);
      \fill[gray] (\x,\y) circle (1.2pt);
    }
    \node at (2.9,1.75) {\small \(t_0\)};
    \node at (2.9,1.25) {\small \(t_1\)};
    \node at (2.9,0.75) {\small \(t_2\)};
    \node at (1.75,2.35) {\small 一个时刻交出一整条剖面};
  \end{scope}
\end{tikzpicture}
\end{center}

\emph{左边是 ODE 的解，右边是 PDE 的解在三个时刻的快照；时间推着整条剖面变形。}

\subsection{变化率由邻域的形状决定，不由本点的取值决定}

ODE 的右端只看当前状态：\(x'=f(x)\)，位置决定速度。方程里没有``邻居''这个词，因为状态空间里的一个点不与别的点相邻。

场就不同了。杆上某一小段是升温还是降温，取决于左边流进来的热量与右边流出去的热量之差，而两侧的热流各由那一侧的斜率给出。斜率之差就是二阶导数。于是某点的走向由该点附近温度曲线的凹凸决定，与该点温度本身有多高无关——一块烧到八百度的金属，若周围同样是八百度，它一动不动。这是从 ODE 跨到 PDE 时最需要重装的直觉。

\subsection{热方程的符号已经把行为写完了}

把上面的收支算清楚。取一小段 \([x,x+\Delta x]\)，左端流入的热流正比于 \(-\partial_xu(t,x)\)，负号来自热量沿温度下降的方向走；右端流出的正比于 \(-\partial_xu(t,x+\Delta x)\)。流入减流出等于这一段内能的积累速率，后者正比于 \(\partial_tu\) 乘以段长。两边同除 \(\Delta x\) 再令 \(\Delta x\to0\)：

\[
\partial_t u=\kappa\,\partial_{xx}u,
\qquad \kappa>0 .
\]

现在只读符号。若某处温度曲线上凸（\(\partial_{xx}u<0\)，该点比左右邻居都热），右端为负，它降温；若下凹（\(\partial_{xx}u>0\)，比邻居都冷），它升温。回到上面那张图右侧：\(t_0\) 时刻峰顶那个灰点处曲线明显上凸，所以峰要往下掉；两侧山脚处曲线下凹，所以肩要往上抬。\(t_1\)、\(t_2\) 两条剖面正是这样从 \(t_0\) 那条压出来的。方程里没有谁规定最终温度是多少，局部的凹凸自己把场推向平坦。

\subsection{差一阶时间导数，弦就从抹平变成来回跑}

把铁杆换成一根拉紧的弦，\(u(t,x)\) 记高度。一小段弦受两侧张力的竖直分量牵拉，合力正比于局部弯曲程度，仍是 \(\partial_{xx}u\)。但牛顿第二定律说的是合力等于质量乘加速度，所以左端出现的是二阶时间导数：

\[
\partial_{tt}u=c^2\,\partial_{xx}u,
\qquad c>0 .
\]

两个方程的右端一模一样，左端差了一阶导。对热方程来说，凸起处的曲率给出一个向下的\emph{速度}，峰一路降到平为止；对波动方程来说，同样的曲率给出一个向下的\emph{加速度}，峰降到平的时候速度最大，于是冲过去，凹到另一侧再被拉回来。一个耗散，一个守恒；一个把高低起伏磨掉，一个让它反复来回。下一节会把这种差别摊开讲。

熟悉的谐振子藏在这里面。把弦的形状按空间频率展开，每个频率分量的振幅满足 \(a''+\omega_k^2a=0\)，正是\hyperref[sec:15-Differential-Equations-and-Dynamical-Systems/01-ODE-Theory/01]{``从增长、衰减与振动说起''}里那个方程。波动方程可以读成无穷多个谐振子被空间结构串在一起。

\subsection{稳态是把时间导数抹掉之后剩下的东西}

热方程一直在磨平起伏。让它跑足够久，起伏磨没了，场不再变化，\(\partial_tu=0\)。代回去：

\[
\partial_{xx}u=0 .
\]

一维稳态温度只能是线性函数，两端的温度值一定，中间就是一条直线。推广到 \(d\) 维空间，得到 Laplace 方程

\[
\Delta u=0,
\qquad
\Delta u=\sum_{i=1}^{d}\partial_{x_ix_i}u ,
\]

其解称为调和函数。方程里没有时间，它描述的是空间上已经停下来的平衡：稳态温度场、静电势、不可压流体的速度势。若每点还有本地的生成或消耗，就得到 Poisson 方程 \(-\Delta u=f\)。\hyperref[sec:08-Numerical-Analysis/06-ODE/04]{``边值问题：从打靶到全局离散''}处理的那根两端定温、内部有热源的金属棒，正是它的一维版本。

\subsection{边界条件是问题的一部分，不是附加说明}

ODE 给一个初值就够了：\(x(0)=x_0\) 加上 \(x'=f(x)\)，往后的轨迹全部确定。场不行。初始条件只管住了时间方向的起点，空间方向还有两个（高维则是一整圈）边缘悬在那里，那里发生什么方程管不着。

一杯刚倒好的咖啡：倒好那一瞬间杯内的温度分布 \(u(0,\cdot)=u_0\) 是初始条件；杯壁如何与外界打交道则要另说。贴着冰袋的那一侧温度被钉住，写作 \(u=0\)，称 Dirichlet 条件；裹了隔热层的那一侧没有热流进出，写作 \(\partial_nu=0\)，称 Neumann 条件；敞开在空气里的那一侧按温差散热，写作 \(\partial_nu+\beta u=g\)，称 Robin 条件。同一个方程配上不同的边界条件，解可以完全不像。

数一数就知道少了会怎样。稳态方程 \(\partial_{xx}u=0\) 的解空间是全体一次函数，两个自由参数，恰好要两个端点值才能钉死。少给一个，得到一族解；给多了且彼此矛盾，一个解也没有。波动方程因为时间导数是二阶，除了初始形状还要给初始速度 \(\partial_tu(0,\cdot)\)，告诉弦松手那一刻正往哪个方向走。至于把条件加在未来时刻——要求弦在 \(t=T\) 停成指定形状——这类问题会立刻变得极其脆弱，第三节专门算这笔账。

方程管局部怎么演化，初值管从哪里出发，边界管在边缘处如何与外界交换。三样凑齐才谈得上唯一解。

\subsection{线性让你可以把初始条件拆开，各自演化再叠回去}

热、波动、Laplace 三条方程都是线性的：\(u\) 与 \(v\) 是解，则 \(\alpha u+\beta v\) 也是解。这份性质带来的好处很实在——复杂的初始分布不必整块对付，可以先拆成简单成分，让每个成分单独演化，末了把结果叠加。

拆成什么成分才划算？拆成空间算子的特征函数。若 \(\phi_k\) 满足 \(\partial_{xx}\phi_k=-\lambda_k\phi_k\) 并符合边界条件，把 \(u\) 沿 \(\set{\phi_k}\) 展开后，每个系数只满足一个标量 ODE，彼此互不打扰。一个 PDE 就此散成一族独立的 ODE。这条路子叫分离变量，第四节走完全程；它与\hyperref[sec:08-Numerical-Analysis/08-Fourier-and-Spectral-Methods/01]{``从 Fourier 级数到连续变换''}中``微分变乘法''说的是同一件事。

非线性 PDE 把这份便利收了回去。Navier--Stokes、Burgers、某些反应扩散方程里，成分之间互相纠缠，拆开后各自演化再叠加得到的东西不再是解。分析和数值的难度都要跳一个量级。本单元守在线性范围内，把三类原型的脾气摸清楚。

下一节换一种看法：不推导，只观察。热方程怎样把一切磨平，波动方程怎样让波形不停地跑，Laplace 方程怎样产出最光滑的函数——这三种形态一旦认得出，拿到一个陌生方程时就知道该期待什么、该防备什么。
